\part{Hamiltonian quantization and relative centres}
The Poisson bracket on a symplectic plane gives a Lie algebra of
Hamiltonians.  Its enveloping algebra quantizes functions on the dual
Lie algebra; its Hochschild cochains describe the relative derived
centre of that quantization.  Write $\hbarq$ for the formal quantization
parameter.  The first comparison in this part is
therefore between two complexes with the same coefficient parameter:
Chevalley cochains with differential $\hbarq d_{\CE}$ and relative
Hochschild cochains of a Rees enveloping algebra.

Finite pointed jets make this comparison explicit.  They also separate
three geometric constructions: a degree-one symplectic target for a
two-dimensional sigma model, a finite odd cotangent space with a quantum
master equation, and the symplectic plane on which Moyal multiplication
acts.  Their coordinates, degrees, and quantizations will be specified
before they are compared.

Two further passages require additional maps.  Formal Hamiltonians form
an inverse limit, whereas their continuous dual is a direct limit.
Matrix Hamiltonians form a necklace Lie algebra, whose abelianization
maps onto commutative Hamiltonians.  Neither passage supplies the
continuous centre comparison by itself.  The positive finite
constructions, a binary matrix homotopy, and the precise comparison
data will give the local models for the global question.

\chapter{Hamiltonian jets and cotangent fields}
\section{Hamiltonians and the Poisson bracket}
Let $V=\C^2$ with coordinates $(z_1,z_2)$ and symplectic form
\[
 \omega=\dd z_1\wedge\dd z_2.
\]
Let $\mathfrak m=(z_1,z_2)$.  The pointed formal Hamiltonian Lie algebra is
\[
 \h=\mathfrak m^2\subset\C\llbracket z_1,z_2\rrbracket,
 \qquad
 \{f,g\}=\partial_{z_1}f\,\partial_{z_2}g-
          \partial_{z_2}f\,\partial_{z_1}g.
\]
It consists of Hamiltonians whose vector fields fix the chosen Darboux
point.  For $p,q\ge0$, the complete filtration
$F^p\h=\mathfrak m^{p+2}$ obeys
\[
 \{F^p\h,F^q\h\}\subset F^{p+q}\h.
\]
Write $R_r=\Sym^r(V^\vee)$ for the homogeneous polynomials of degree $r$.  The bracket obeys
\[
 \{R_r,R_s\}\subset R_{r+s-2}.
\]
For $d\ge2$ define the finite jet quotient
\[
 \h_{\le d}=\bigoplus_{r=2}^{d}R_r,
\]
where every bracket component of degree greater than $d$ is set equal to zero.

\begin{theorem}[Finite Hamiltonian jet algebra]
The truncated bracket makes $\h_{\le d}$ a finite-dimensional Lie algebra.  Its dimension is
\[
 \dim\h_{\le d}=\frac{d(d+3)}2-2.
\]
The quadratic part is
\[
 R_2\cong\mathfrak{sl}_2,
\]
and each $R_r$ is the irreducible $\mathfrak{sl}_2$-module $\Sym^r(\C^2)$.
\end{theorem}
\begin{proof}
The Jacobiator of the polynomial Poisson bracket vanishes before truncation.  Projection to degrees at most $d$ preserves this identity because the discarded terms form a Lie ideal.  The dimension is
\[
 \sum_{r=2}^{d}(r+1)
 =\frac{d(d+3)}2-2.
\]
Quadratic Hamiltonians act linearly on $V$ and preserve $\omega$; this identifies $R_2$ with $\mathfrak{sp}(V)=\mathfrak{sl}_2$.  The standard action on homogeneous polynomials gives the final statement.
\end{proof}

For monomials $z^\alpha=z_1^az_2^b$ and
$z^\beta=z_1^cz_2^d$, the bracket is
\[
 \{z^\alpha,z^\beta\}
 =\det(\alpha,\beta)z^{\alpha+\beta-\epsilon},
 \qquad \epsilon=(1,1).
\]
The coefficient of the Jacobiator on three monomials is
\[
 \sum_{\mathrm{cyc}}
 \det(\alpha,\beta)\det(\alpha+\beta-\epsilon,\gamma)=0.
\]
Expanding the second determinant cancels its quadratic terms in pairs.
The remaining terms vanish by
$\det(\alpha,\beta)\gamma+
\det(\beta,\gamma)\alpha+
\det(\gamma,\alpha)\beta=0$ in dimension two.
This is the coordinate form of the identity
$[\partial_{z_1}\wedge\partial_{z_2},
\partial_{z_1}\wedge\partial_{z_2}]_{\mathrm{Sch}}=0$.

The affine Hamiltonian Lie algebra
$\h^{\mathrm{aff}}=\C\llbracket z_1,z_2\rrbracket/\C$
also contains linear Hamiltonians, which generate translations.
It does not have the same finite jet quotients: for $d\ge2$, the
bracket of $z_1^{d+1}$ with $z_2$ is $(d+1)z_1^d$, so the discarded
higher-degree subspace is no longer an ideal.

\section{Unimodularity}
\begin{theorem}[Unimodularity of every finite jet]
For every $d\ge2$ and every $x\in\h_{\le d}$,
\[
 \Tr_{\h_{\le d}}(\ad_x)=0.
\]
\end{theorem}
\begin{proof}
A homogeneous Hamiltonian of degree greater than two raises polynomial degree.  Its adjoint action is strictly upper triangular in the decomposition $\bigoplus_{r=2}^{d}R_r$, hence has trace zero.  A quadratic Hamiltonian belongs to $\mathfrak{sl}_2$.  Its action on each finite-dimensional $\mathfrak{sl}_2$-module $R_r$ has trace zero.  Linearity gives the result.
\end{proof}


\section{The degree-one symplectic target}
Use cohomological grading, with $(M[k])^i=M^{i+k}$.
Put $L=\h_{\le d}$ in degree zero and define the semidirect dg Lie algebra
\[
 \g_d=L\ltimes L^\vee[-1].
\]
The dual summand has degree $1$, is abelian, and carries the coadjoint
action $(\ad_x^*\lambda)(y)=-\lambda([x,y])$.
Thus $\g_d[1]=L[1]\oplus L^\vee$ is the degree-one symplectic
space $T^*[1]L^\vee$.  Its functions have even coordinates $x_a$ on
$L^\vee$ and odd coordinates $c^a$ of degree $1$ in the cotangent
direction.  These $x_a$ are linear functions corresponding to a basis
$e_a$ of $L$; write $[e_a,e_b]=f^c{}_{ab}e_c$.

\begin{proposition}[Chevalley--Poisson identification]
There is a canonical isomorphism of commutative dg algebras
\[
 C^*_{\CE}(\g_d)
 \cong C^*(L,\Sym L)
 =\bigl(\Sym L\tensor\Lambda L^\vee,d_{\CE}\bigr).
\]
Under $\Sym L=\Ocal(L^\vee)$, this is the polynomial polyvector
complex for the linear Poisson tensor on $L^\vee$.
\end{proposition}
\begin{proof}
Chevalley functions are $\Sym(\g_d^\vee[-1])$.
The dual of $L^\vee[-1]$, followed by the shift $[-1]$, is $L$ in
degree zero.  The other summand contributes $L^\vee[-1]$ in degree
$1$.  The Lie bracket and coadjoint action give
\begin{equation}\label{eq:vol2-ce-coordinate}
 d_{\CE}c^c=-\frac12f^c{}_{ab}c^ac^b,
 \qquad
 d_{\CE}x_a=f^c{}_{ba}c^b x_c.
\end{equation}
The Schouten differential of the linear Poisson tensor has these same
values on polynomial coordinates and constant vector fields.
Both differentials are derivations, so the identifications agree on
the entire algebra.
\end{proof}

In particular, for a polynomial $F$ one has
\[
 d_{\CE}F=c^a\{x_a,F\}_{\mathrm{lin}},
 \qquad \{x_a,x_b\}_{\mathrm{lin}}=f^c{}_{ab}x_c.
\]
The identity $d_{\CE}^2=0$ is the Poisson Jacobi identity.
With the degree-$-1$ Poisson bracket $\{c^a,x_b\}=\delta^a_b$,
the Hamiltonian
\[
 \Theta=\frac12x_cf^c{}_{ab}c^ac^b
\]
has degree $2$ and generates this differential.
It is the Hamiltonian of the degree-one symplectic target.

\section{The finite BV space}
The finite BV construction instead uses
\[
 \mathcal E_L=T^*[-1](L[1])=L[1]\oplus L^\vee[-2].
\]
Its linear coordinates are $c^a$ of degree $1$ and $b_a$ of degree
$-2$.  Choose the odd Poisson bracket $\{b_a,c^b\}_{\mathrm{BV}}
=\delta_a^b$ and the constant Berezinian, with left odd derivatives
in the Laplacian
\[
 \Delta=\sum_a\frac{\partial}{\partial b_a}
                    \frac{\partial}{\partial c^a}.
\]
The degree-zero cotangent action is
\[
 S_L=-\frac12b_cf^c{}_{ab}c^ac^b.
\]
Its Hamiltonian vector field on the $c$ coordinates is the Chevalley
differential in \eqref{eq:vol2-ce-coordinate}.

\begin{theorem}[Finite quantum master equation]
For $L=\h_{\le d}$ the constant Berezinian gives
\[
 \frac12\{S_L,S_L\}_{\mathrm{BV}}=0,
 \qquad
 \Delta S_L=\sum_b\Tr_L(\ad e_b)c^b=0.
\]
Consequently $S_L$ satisfies the finite quantum master equation
\[
 \frac12\{S_L,S_L\}_{\mathrm{BV}}+\hbarq\Delta S_L=0.
\]
\end{theorem}
\begin{proof}
The classical term is the cotangent Hamiltonian for the square of the
Chevalley differential, hence vanishes by Jacobi.
Left differentiation of the two odd variables gives
\[
 \Delta S_L=-\sum_{a,b}f^a{}_{ab}c^b
            =\sum_b\Tr_L(\ad e_b)c^b.
\]
Unimodularity of the finite Hamiltonian jet algebra gives the second
vanishing.  The chosen Berezinian is translation invariant; multiplying
it by a nonzero constant does not change $\Delta$.
\end{proof}

The two cotangent constructions have different degrees and functions.
The even variables $x_a$ in the Poisson complex have degree zero;
the finite BV antifields $b_a$ have degree $-2$.
The finite master equation is an algebraic statement on
$\mathcal E_L$.  A field theory requires a source manifold, its field
complex, and boundary conditions.

\section{The Poisson sigma model and its boundary product}
Let $\Sigma$ be an oriented disk and use the degree-one symplectic
target $T^*[1]L^\vee$.  The graded mapping space
$\operatorname{Map}(T[1]\Sigma,T^*[1]L^\vee)$ has a degree-$-1$
symplectic form.  Its degree-zero classical fields are a map
$X:\Sigma\to L^\vee$ and an $L$-valued one-form $\eta$.
For the linear Poisson tensor $\pi$ on $L^\vee$, the classical action is
\[
 S_\Sigma(X,\eta)
 =\int_\Sigma\left(\eta_a\wedge\dd X^a
        +\frac12\pi^{ab}(X)\eta_a\wedge\eta_b\right).
\]
Impose $\eta|_{T\partial\Sigma}=0$.
Integration by parts identifies the action with
\[
 \int_\Sigma\left\langle X,\dd\eta+
                              \frac12[\eta,\eta]\right\rangle.
\]
The boundary term vanishes under the stated condition.
This is two-dimensional BF theory for $L$; the full mapping-space fields add
the ghosts and antifields required by its gauge symmetry.

For boundary insertions at two ordered points and a fixed value of
$X$ at a third point, the formal perturbation expansion gives the
Kontsevich graph product.  The propagator and the finite counterterm
which removes tadpoles are described in
\cite{CattaneoFelderPSM}*{Sections~3.7--3.8}.
This is a perturbative interpretation of the product.  Its associativity
and the following algebraic identification use the formality theorem.

\begin{theorem}[Linear Poisson boundary algebra]
Let $R=\C\llbracket\hbarq\rrbracket$ and let $\star_L$ be the
Kontsevich product for $\hbarq\pi$ on $\Sym L\llbracket\hbarq\rrbracket$.
There is an $R$-algebra isomorphism
\[
 (\Sym L\llbracket\hbarq\rrbracket,\star_L)
 \xrightarrow{\sim}
 U_\hbar(L)
 =T(L)\llbracket\hbarq\rrbracket/
   \overline{\angles{xy-yx-\hbarq[x,y]}}.
\]
It fixes the linear generators.  The associative algebra on the right
defines a locally constant factorization algebra on the oriented line,
whose product for adjacent intervals is the enveloping multiplication.
\end{theorem}
\begin{proof}
For linear functions the star commutator is exactly
$\hbarq[x,y]$.  The universal property gives the homomorphism from
the Rees enveloping algebra to the star algebra.
Its reduction modulo $\hbarq$ is the identity on $\Sym L$.
PBW and completeness lift this to an isomorphism; the inverse is the
displayed map.  This is the linear Poisson identification of
\cite{KontsevichDQ}*{Section~8.2}.
The usual associative factorization envelope on a line assigns the
algebra to an interval and its ordered multiplication to disjoint
subintervals.  Associativity supplies its gluing maps.
\end{proof}

PBW symmetrization is the underlying $R$-module isomorphism.
It also transports the enveloping product to the Gutt product.
The equivalence with the Kontsevich product is an algebra map and
contains its own ordering correction.  Identifying an analytic or
renormalized field-theoretic boundary observable algebra with this
algebraic factorization model requires the corresponding observable map.

\chapter{Relative centres and affine formality}
\section{The coefficient ring and the Chevalley differential}
The Rees relation scales both the Lie bracket and its adjoint action.
Its Chevalley differential must retain the same parameter.
For example, take the quadratic jet $\flie=\mathfrak{sl}_2$ with
$[h,e]=2e$.  Modulo $\hbarq$, the enveloping algebra is commutative,
so the Hochschild differential of $e$ vanishes.
The unscaled Chevalley differential instead satisfies
$(d_{\CE}e)(h)=2e$.

Let $\flie$ be a finite-dimensional complex Lie algebra.
For $N\ge0$, put
\[
 R_N=\C[\hbarq]/(\hbarq^{N+1}),\qquad
 \flie_N=R_N\tensor_\C\flie,\qquad
 [x,y]_N=\hbarq[x,y],
\]
and
\[
 A_N=U_{R_N}(\flie_N)
 =T_{R_N}(\flie_N)/(xy-yx-\hbarq[x,y]).
\]
Here $N$ denotes the order in $\hbarq$, independently of a Chan--Paton rank.
Write $\bar A_N$ for the augmentation ideal and $S=\Sym_\C\flie$.
PBW symmetrization is the $R_N$-module isomorphism
\[
 \operatorname{sym}_N:S\tensor R_N\longrightarrow A_N,
 \qquad
 x_1\cdots x_r\longmapsto
 \frac1{r!}\sum_{\sigma\in S_r}x_{\sigma(1)}\cdots x_{\sigma(r)}.
\]
The source has its ordinary symmetric product.
Use normalized relative Hochschild cochains
\[
 C^n_{\HH}(A_N/R_N,A_N)
 =\operatorname{Hom}_{R_N}(\bar A_N^{\tensor_{R_N}n},A_N)
\]
with differential
\[
\begin{split}
 (bF)(a_1,\ldots,a_{n+1})
 ={}&a_1F(a_2,\ldots,a_{n+1})\\
 &+\sum_{i=1}^n(-1)^iF(a_1,\ldots,a_i a_{i+1},\ldots,a_{n+1})\\
 &+(-1)^{n+1}F(a_1,\ldots,a_n)a_{n+1}.
\end{split}
\]
Let $d_{\CE}$ act on
$C^n(\flie,S)=\operatorname{Hom}_\C(\Lambda^n\flie,S)$
by the adjoint Chevalley differential.

\begin{theorem}[Relative Rees comparison]\label{thm:vol2-relative-rees}
Antisymmetrization defines a quasi-isomorphism of $R_N$-complexes
\[
 \operatorname{Alt}_N:
 C^\bullet_{\HH}(A_N/R_N,A_N)
 \longrightarrow
 \bigl(C^\bullet(\flie,S)\tensor R_N,\hbarq d_{\CE}\bigr),
\]
given by
\[
 (\operatorname{Alt}_NF)(x_1,\ldots,x_n)
 =\operatorname{sym}_N^{-1}
 \sum_{\sigma\in S_n}\operatorname{sgn}(\sigma)
 F(x_{\sigma(1)},\ldots,x_{\sigma(n)}).
\]
These maps commute with coefficient reduction in $N$.
\end{theorem}
\begin{proof}
Write $A=A_N$, $R=R_N$, and $A^e=A\tensor_R A^{\op}$.
In homological degree $n$, set
\[
 K_n=A\tensor_R\Lambda_R^n\flie_N\tensor_R A.
\]
Its augmentation is multiplication.
For $X=x_1\wedge\cdots\wedge x_n$, let $X_{\widehat i}$ and
$X_{\widehat i\widehat j}$ omit the indicated entries without reordering.
Define
\[
\begin{split}
 \partial(a\tensor X\tensor b)
 ={}&\sum_i(-1)^{i+1}
 \bigl(ax_i\tensor X_{\widehat i}\tensor b
       -a\tensor X_{\widehat i}\tensor x_i b\bigr)\\
 &+\hbarq\sum_{i<j}(-1)^{i+j}
 a\tensor[x_i,x_j]\wedge X_{\widehat i\widehat j}\tensor b.
\end{split}
\]
The normalized relative bimodule bar has terms
$\mathcal B_n=A\tensor_R\bar A^{\tensor_R n}\tensor_R A$.
Define
\[
 \iota_n(a\tensor X\tensor b)
 =\sum_{\sigma\in S_n}\operatorname{sgn}(\sigma)
       a[x_{\sigma(1)}|\cdots|x_{\sigma(n)}]b.
\]
The outer bar faces give the first sum in $\partial$.
Pair adjacent transpositions in the internal faces and use
$xy-yx=\hbarq[x,y]$ to obtain the second sum.
For two inputs this reads
\[
\begin{split}
 b\bigl(1[x|y]1-1[y|x]1\bigr)
 ={}&x[y]1-y[x]1+1[x]y-1[y]x\\
 &-\hbarq\,1[[x,y]]1.
\end{split}
\]
Thus $b\iota=\iota\partial$.
PBW makes $\flie_N$ a free summand of $\bar A$.
Since every factorial is invertible in $R$, $\iota$ is injective.
The identity $b^2=0$ then gives $\partial^2=0$.

Give $K_n$ filtration degree equal to its outer PBW degrees plus $n$.
The associated graded is the diagonal Koszul complex of
$\Sym_R(\flie_N)$.
Its differential uses differences of left and right polynomial generators.
A linear change of variables makes these independent polynomial variables.
They form a regular sequence even when $R$ has nilpotents.
For an augmented cycle, lift a primitive of its highest-degree symbol
and subtract its differential.
The filtration degree decreases, so this procedure terminates.
Hence $K_\bullet\to A$ is a free bimodule resolution.
The relative bar is also free over $A^e$, with its usual unit-insertion contraction.
The map $\iota$ lifts the identity of $A$.
Projectivity constructs its chain homotopy inverse degree by degree.
Applying $\operatorname{Hom}_{A^e}(-,A)$ gives a cochain homotopy equivalence.

The induced differential on
$\operatorname{Hom}_R(\Lambda_R^n\flie_N,A)$ is
\[
\begin{split}
 (\delta_N\phi)(x_1,\ldots,x_{n+1})
 ={}&\sum_i(-1)^{i+1}
 [x_i,\phi(x_1,\ldots,\widehat{x_i},\ldots,x_{n+1})]_A\\
 &+\hbarq\sum_{i<j}(-1)^{i+j}
 \phi([x_i,x_j],x_1,\ldots,\widehat{x_i},
                         \ldots,\widehat{x_j},\ldots,x_{n+1}).
\end{split}
\]
Commuting $x$ past the factors of a symmetrized monomial gives
\[
 [x,\operatorname{sym}_N(f)]_A
 =\hbarq\operatorname{sym}_N(\operatorname{ad}_x f).
\]
Both sums therefore become $\hbarq$ times the ordinary Chevalley differential.
Precomposition with $\iota$ gives $\operatorname{Alt}_N$.
All formulas commute with coefficient reduction.
\end{proof}

\section{The formal coefficient limit}
Put $R=\C\llbracket\hbarq\rrbracket$ and $A=\varprojlim_N A_N$.
PBW identifies its underlying module with $S\llbracket\hbarq\rrbracket$.
Each coefficient of $\hbarq$ is a polynomial.
Define
\[
 A^{\widehat\tensor_R n}=\varprojlim_N A_N^{\tensor_{R_N}n},
 \qquad
 C^\bullet_{\HH,\mathrm{cont}}(A/R,A)
 =\varprojlim_N C^\bullet_{\HH}(A_N/R_N,A_N).
\]
Equivalently, these are continuous $R$-multilinear normalized cochains
on the completed relative tensors.

\begin{corollary}[Formal relative Rees comparison]\label{cor:vol2-formal-relative-rees}
The compatible antisymmetrizations induce a quasi-isomorphism
\[
 C^\bullet_{\HH,\mathrm{cont}}(A/R,A)
 \longrightarrow
 \bigl(C^\bullet(\flie,S)\llbracket\hbarq\rrbracket,
                    \hbarq d_{\CE}\bigr).
\]
\end{corollary}
\begin{proof}
Coefficient reduction is surjective in each cochain degree.
For Hochschild cochains, lift outputs on the PBW input basis.
The Chevalley exterior input has finite rank, so its outputs also lift.
The comparison cones are acyclic and their transition maps are degreewise surjective.
The long exact sequence shows that every transition kernel is acyclic.
Choose a primitive of a compatible cycle at level zero.
Lift it one level and correct its differential by a primitive in the transition kernel.
Iteration gives compatible primitives at every order.
The inverse-limit cone is therefore acyclic.
\end{proof}

The coefficient ring is part of this result.
Absolute $\C$-linear cochains also detect variations of the coefficient parameter.
For $\flie=0$, the continuous absolute derivation $\partial_{\hbarq}$
on $R$ has no relative counterpart.
For nonabelian $\flie$, the associated graded for the $\hbarq$ filtration
has zero differential, whereas $d_{\CE}$ need not vanish.

\section{Linear Poisson formality}
For $\flie=\h_{\le d}$, put $A_d=U_\hbar(\h_{\le d})$ and
\[
 \mathcal P_d=
 \bigl(C^*(\h_{\le d},\Sym\h_{\le d})
       \llbracket\hbarq\rrbracket,\hbarq d_{\CE}\bigr).
\]
This is the polynomial polyvector complex on $\h_{\le d}^*$ with
differential $[\hbarq\pi_d,-]$ for the linear Poisson tensor $\pi_d$.
The relative comparison gives $C^\bullet_{\HH,\mathrm{cont}}(A_d/R,A_d)
\to\mathcal P_d$ as a map of complexes.
Higher operations require a formality map with specified operadic structure.

Fix Kontsevich's formality morphism $\mathcal U$ on the affine space
$\h_{\le d}^*$.
Its twist by $\hbarq\pi_d$ defines a star product $\star_d$ and an
$\Linf$ quasi-isomorphism on the degree-shifted complexes.
Let $T_d:(\Sym\h_{\le d}\llbracket\hbarq\rrbracket,\star_d)
\to A_d$ be the algebra identification fixing the linear generators.
This identification includes the equivalence between the Kontsevich and PBW products.
PBW symmetrization alone identifies the underlying modules.
The first Taylor coefficient of the twisted comparison is
\[
 \Phi_d(v)=(T_d)_*
 \sum_{r\ge0}\frac1{r!}
 \mathcal U_{r+1}
 (\underbrace{\hbarq\pi_d,\ldots,\hbarq\pi_d}_{r},v).
\]
It has source $\mathcal P_d$ and target
$C^\bullet_{\HH,\mathrm{cont}}(A_d/R,A_d)$.
At each order in $\hbarq$, polynomial inputs have polynomial outputs.
Its reduction modulo $\hbarq$ is the Hochschild--Kostant--Rosenberg map,
with the normalized antisymmetrization convention.
That map is a quasi-isomorphism for the polynomial algebra.
The complete $\hbarq$ filtration then proves that $\Phi_d$ is a quasi-isomorphism.
The higher Taylor coefficients supply the $\Linf$ identities after shifting by $[1]$.
These statements use the affine formality construction
\cite{KontsevichDQ}*{Sections~6, 8.2, and~8.3}.

On invariant polynomials, the standard Duflo normalization gives
\[
 \operatorname{Duf}_\hbar
 =\operatorname{sym}_\hbar\circ j^{1/2}(\hbarq\partial),
 \qquad
 j^{1/2}(x)=
 \det_{\h_{\le d}}^{1/2}
 \left(\frac{\sinh(\operatorname{ad}_x/2)}
 {\operatorname{ad}_x/2}\right).
\]
It is an algebra isomorphism from
$(\Sym\h_{\le d})^{\h_{\le d}}\llbracket\hbarq\rrbracket$
to $Z(A_d)$, by the Duflo theorem
\cite{KontsevichDQ}*{Section~8.3.4}.
The square root has constant term one.
This assertion concerns invariant degree-zero cohomology.
It does not identify antisymmetrization with a homotopy-$\Etwo$ map.

\section{The quadratic Duflo correction}
First set the coefficient parameter to $1$ in the polynomial Rees algebra.
For the resulting ordinary $\mathfrak{sl}_2$ enveloping algebra, take generators
\[
 [h,e]=2e,
 \qquad
 [h,f]=-2f,
 \qquad
 [e,f]=h.
\]
The invariant polynomial and the quadratic Casimir are
\[
 q=2ef+\frac12h^2,
 \qquad
 \Omega=ef+fe+\frac12h^2.
\]

\begin{calculation}[The zero-point shift]
The Duflo map gives
\[
 \operatorname{Duf}(q)=\Omega+\frac12.
\]
Let $V_n$ be the irreducible highest-weight-$n$ module, with $n\ge0$.
Then
\[
 \left(\Omega+\frac12\right)\big|_{V_n}
 =\frac{(n+1)^2}{2}\id.
\]
\end{calculation}
\begin{proof}
Write $x=ue+vf+wh$.  Direct adjoint matrices give
$\Tr(\operatorname{ad}_x^2)=8uv+8w^2$.
Thus the quadratic term of $j^{1/2}(\partial)$ is
$(\partial_e\partial_f+\partial_h^2)/6$.
Applied to $q=2ef+h^2/2$, it gives $(2+1)/6=1/2$.
Higher differential terms annihilate this quadratic polynomial.
Symmetrization sends $q$ to $\Omega$.
On a highest-weight vector, $ef$ acts by $n$, $fe$ by zero,
and $h^2/2$ by $n^2/2$.
Centrality gives the eigenvalue on the entire irreducible module.
\end{proof}

Restoring the Rees parameter gives $\operatorname{Duf}_\hbar(q)=\Omega+\hbarq^2/2$.  The polynomial calculation permits specialization at $1$; arbitrary formal power series do not admit that specialization.  Higher wheels assemble into $j^{1/2}$ on higher-degree invariants.


For the quadratic jet, the invariant polynomial $q=2ef+h^2/2$ is
$d_{\CE}$-closed.  Its Duflo image is a degree-zero central class.
The oscillator quotient below sends that class to $\hbarq^2/8$;
this scalar follows from the quotient's own Casimir relation.

\section{Hochschild braces}
For Hochschild cochains $F\in C^p(A,A)$ and $G_i\in C^{q_i}(A,A)$, the brace
\[
 F\{G_1,\ldots,G_r\}
\]
is the signed sum over ordered insertions of the $G_i$ into the inputs
of $F$.  For an algebra in degree zero, an insertion has sign
$(-1)^{\sum_i(q_i-1)j_i}$, where $j_i$ is the number of input
positions preceding $G_i$ in the expanded expression.
The brace identities state that a nested insertion is the sum over
all ways of distributing the inner insertions among the outer slots.
Together with the cup product and unit, these operations give the
Hochschild form of the Deligne structure
\cite{McClureSmith}*{Theorem~1.1}.
The operadic construction in
\cite{KontsevichSoibelmanDeligne}*{Section~7} gives the corresponding
structure for $A_\infty$ algebras.  Here the relative operations are
taken at each coefficient precision and then passed to the complete
limit; every fixed-arity insertion is continuous.

The relative Hochschild cochains carry this brace structure.
Their canonical coherent evaluation action on $A_d$ gives the
closed--open pair
$(C^\bullet_{\HH,\mathrm{cont}}(A_d/R,A_d),A_d)$.
The closed algebra is the relative derived centre.
A homotopy-$\Etwo$ refinement means a morphism or zigzag
\[
 \widetilde\Phi_d:\mathcal P_d
 \simeq C^\bullet_{\HH,\mathrm{cont}}(A_d/R,A_d)
\]
in complete $R$-linear homotopy-$\Etwo$ algebras,
with the chosen Gerstenhaber structure on $\mathcal P_d$.
A compatible Swiss-cheese refinement also specifies the mixed operations
on the pair and their coherence homotopies.

\begin{corollary}[Open--closed structure with a refinement]\label{cor:vol2-open-closed-refinement}
Given these refinements, the pair $(\mathcal P_d,A_d)$ has the transported
homotopy Swiss-cheese action.
Its open product is the enveloping product.
The specified closed and mixed Taylor coefficients determine its
closed bracket and boundary action.
\end{corollary}
\begin{proof}
Transport the two-colored operations through the supplied equivalence
of complete operadic algebras.
The coherence homotopies are part of this equivalence.
\end{proof}

\begin{theorem}[Brace realization with a refinement]
For $A=A_d$, supply the homotopy-$\Etwo$ and mixed refinements of
Corollary~\ref{cor:vol2-open-closed-refinement}.
They identify the selected homotopy Gerstenhaber operations on
$\mathcal P_d$ with the relative Hochschild centre operations.
The chosen mixed refinement also identifies their action on the boundary algebra.
\end{theorem}
\begin{proof}
The relative Hochschild complex has its brace operations.
The supplied operadic equivalence preserves their coherent homotopies.
An equivalence of complexes alone would not imply this conclusion.
\end{proof}

\section{The first mixed identities}
A relative Hochschild $1$-cocycle $D$ of an algebra concentrated in degree zero
is an $R$-linear derivation:
\[
 D(ab)=D(a)b+aD(b).
\]
For two such cocycles $D,E$, their Gerstenhaber bracket acts by
\[
 [D,E](a)=D(E(a))-E(D(a)).
\]
For a graded algebra and homogeneous derivations, these formulas become
\[
\begin{split}
 D(ab)&=D(a)b+(-1)^{|D||a|}aD(b),\\
 [D,E](a)&=D(E(a))-(-1)^{|D||E|}E(D(a)).
\end{split}
\]
Higher Hochschild cochains act through multilinear evaluations and braces.
They need not act by derivations or carry a coproduct defining a diagonal action.
A compatible mixed refinement transports these operations, including their homotopies.

\chapter{Finite presentations and Moyal multiplication}
\section{Quadratic Hamiltonians}
Set
\[
 e=\frac12z_1^2,
 \qquad
 f=-\frac12z_2^2,
 \qquad
 h=-z_1z_2.
\]
A direct differentiation gives
\[
 \{h,e\}=2e,
 \qquad
 \{h,f\}=-2f,
 \qquad
 \{e,f\}=h.
\]
Thus $\h_{\le2}=R_2$ is $\mathfrak{sl}_2$ with its standard
Chevalley basis.  Its Rees enveloping algebra is
$U_\hbar(\mathfrak{sl}_2)$; the Moyal realization below is a proper
associative quotient of this algebra.

\section{Cubic Hamiltonians}
Let
\[
 v_j=z_1^{3-j}z_2^j,
 \qquad 0\le j\le3.
\]
Then
\[
 \{h,v_j\}=(3-2j)v_j,
 \qquad
 \{e,v_j\}=jv_{j-1},
 \qquad
 \{f,v_j\}=(3-j)v_{j+1},
\]
with the boundary terms understood as zero.  Hence $R_3$ is the irreducible highest-weight-three module.

\begin{proposition}[The cubic jet algebra]
The degree-three jet algebra is the semidirect product
\[
 \h_{\le3}\cong\mathfrak{sl}_2\ltimes\Sym^3(\C^2),
\]
where the second summand is abelian.
\end{proposition}
\begin{proof}
The bracket of two cubic polynomials has degree four and vanishes in the degree-three quotient.  The bracket with a quadratic polynomial is the displayed $\mathfrak{sl}_2$ action.
\end{proof}

\section{The first nonlinear layer}
Write $w_k=z_1^{4-k}z_2^k$, $0\le k\le4$.  In $\h_{\le4}$ the cubic modes satisfy
\begin{equation}\label{eq:vol2-cubic-quartic}
 \{v_i,v_j\}=3(j-i)w_{i+j-1},
\end{equation}
whenever the index lies between $0$ and $4$; the coefficient makes every remaining case vanish.

\begin{calculation}[Two brackets]
\[
 \{z_1^3,z_2^3\}=9z_1^2z_2^2=9w_2,
 \qquad
 \{z_1^2z_2,z_1z_2^2\}=3z_1^2z_2^2=3w_2.
\]
The two independent cubic pairs reach the same quartic channel with amplitudes in the ratio $3:1$.
\end{calculation}
\begin{proof}
Differentiate each monomial in the two symplectic coordinates and substitute in the Poisson formula.  The first pair gives $3z_1^2\cdot3z_2^2$.  The second gives $2z_1z_2\cdot2z_1z_2-z_1^2\cdot z_2^2=3z_1^2z_2^2$.
\end{proof}

\begin{theorem}[The quartic jet presentation]
The twelve symbols $e,f,h,v_0,\ldots,v_3,w_0,\ldots,w_4$ present
$\h_{\le4}$ with the following relations: the $\mathfrak{sl}_2$
relations, the displayed action on $v_j$,
\eqref{eq:vol2-cubic-quartic}, and
\[
\begin{aligned}
 [h,w_k]&=(4-2k)w_k,&[e,w_k]&=kw_{k-1},&
 [f,w_k]&=(4-k)w_{k+1},\\
 [v_i,w_k]&=0,&[w_k,w_l]&=0.
\end{aligned}
\]
Indices outside their indicated ranges contribute zero.
The seven symbols $e,f,h,v_0,\ldots,v_3$ already generate this Lie algebra.
\end{theorem}
\begin{proof}
Every bracket of two displayed symbols is a linear combination of the
twelve symbols.  Induction on a Lie word therefore shows that the
presented algebra has dimension at most twelve.
The polynomial Poisson bracket, truncated above degree four, satisfies
every relation.  Indeed, a cubic--quartic bracket has degree five and a
quartic--quartic bracket has degree six, so both vanish in this quotient.
The images of the twelve symbols form the monomial basis of
$R_2\oplus R_3\oplus R_4$, which has dimension twelve.
The resulting surjection is consequently an isomorphism.
Finally, \eqref{eq:vol2-cubic-quartic} produces $w_0,w_1,w_2,w_3,w_4$
from the pairs $(v_0,v_1),(v_0,v_2),(v_0,v_3),(v_1,v_3),(v_2,v_3)$,
with nonzero coefficients $3,6,9,6,3$.
\end{proof}

\section{The Moyal product}
For polynomials on $\C^2$, the Moyal product is
\begin{equation}\label{eq:vol2-moyal}
 f\star g
 =\sum_{n\ge0}\frac1{n!}\left(\frac{\hbarq}{2}\right)^n
 \sum_{r=0}^n(-1)^r\binom nr
 (\partial_1^{n-r}\partial_2^rf)
 (\partial_1^r\partial_2^{n-r}g).
\end{equation}
Every polynomial pair gives a finite sum.
Associativity follows directly from the constant Poisson tensor.
On a triple tensor product, the operators
$P_{ij}=\partial_1^{(i)}\partial_2^{(j)}-
\partial_2^{(i)}\partial_1^{(j)}$ commute.  Both parenthesizations
of the product are therefore
\[
 m_3\exp\!\left(\frac{\hbarq}{2}
                   (P_{12}+P_{13}+P_{23})\right)
        (f\tensor g\tensor h),
\]
where $m_3$ multiplies the three polynomial factors.

\begin{theorem}[Quadratic exactness]
For polynomials $f,g$ of degree at most two,
\[
 [f,g]_\star=\hbarq\{f,g\}.
\]
The quadratic generators therefore define a Lie representation of the Rees-scaled $\mathfrak{sl}_2$.  Their associative algebra has an additional Casimir relation.
\end{theorem}
\begin{proof}
The commutator retains the odd powers of the Poisson bidifferential.  Derivatives of order three vanish on quadratics, so the first-order term is the entire commutator.
\end{proof}

\begin{proposition}[The oscillator quotient]\label{prop:vol2-oscillator-quotient}
Let $R_N=\C[\hbarq]/(\hbarq^{N+1})$ and let $W_N=R_N[q,p]$
have the Moyal product for $\{q,p\}=1$.
Let $U_N$ have generators $E,F,H$ with relations
\[
 [H,E]=2\hbarq E,\qquad [H,F]=-2\hbarq F,\qquad [E,F]=\hbarq H.
\]
Set $\Omega=EF+FE+H^2/2$.
The assignments
\[
 E\longmapsto e=q^2/2,\qquad
 F\longmapsto f=-p^2/2,\qquad
 H\longmapsto h=-qp
\]
induce an algebra isomorphism
\[
 U_N/(\Omega+3\hbarq^2/8)
 \xrightarrow{\sim} W_N^{(q,p)\mapsto(-q,-p)}.
\]
These compatible isomorphisms also hold after $\hbarq$-adic completion.
If linear Hamiltonians are retained, the generated associative algebra is the full Weyl algebra.
\end{proposition}
\begin{proof}
Quadratic exactness proves the three commutator relations.
The associative products are
\[
\begin{split}
 e\star f&=-q^2p^2/4-\hbarq qp/2-\hbarq^2/8,\\
 f\star e&=-q^2p^2/4+\hbarq qp/2-\hbarq^2/8,\\
 h\star h&=q^2p^2-\hbarq^2/4.
\end{split}
\]
Thus $\Omega$ maps to $-3\hbarq^2/8$.
Its commutator with each of $E,F,H$ vanishes by the defining relations.
Give these generators PBW degree one.
Give an even monomial $q^ap^b$ degree $(a+b)/2$.
The leading symbol of $\Omega+3\hbarq^2/8$ is
$C=2EF+H^2/2$.
It is a non-zero-divisor in $R_N[E,F,H]$, since $2C$ is monic in $H$.
Centrality and this property give
\[
 \operatorname{gr}\bigl(U_N/(\Omega+3\hbarq^2/8)\bigr)
 =R_N[E,F,H]/(C).
\]
Indeed, every element of the two-sided ideal is a multiple of the central generator.
Its leading symbol is the product of the two leading symbols.
The quotient has basis $E^aH^bF^c$ with $a,c\ge0$ and $b=0,1$.
Their leading symbols map to unit multiples of the distinct monomials
$q^{2a+b}p^{2c+b}$, precisely the even monomials.
Hence the associated-graded map is an isomorphism.
Lifting highest-degree terms proves surjectivity and comparing them proves injectivity.
Every element at fixed $N$ has finite polynomial degree.
The formulas commute with reduction in $N$, so their inverse limit gives the formal isomorphism.
Finally $q,p$ generate $W_N$ and topologically generate its formal limit.
\end{proof}

The primitive coproduct on $U_N$ does not descend to this quotient.
Modulo $\hbarq$, it sends the classical relation to
\[
 \Delta(C)=C\tensor1+1\tensor C
   +2E\tensor F+2F\tensor E+H\tensor H.
\]
The cross-term is nonzero in the tensor square of $\C[E,F,H]/(C)$.
For example, evaluation at $(E,F,H)=(1,0,0)$ in the first factor
and $(0,1,0)$ in the second gives $2$.
Thus the oscillator quotient is an associative algebra, with no coproduct inherited from these primitive generators.

\begin{calculation}[Cubic correction]
For $f=z_1^3$ and $g=z_2^3$,
\[
 f\star g-g\star f
 =9\hbarq z_1^2z_2^2+\frac32\hbarq^3.
\]
The first term quantizes the quartic Hamiltonian channel.  The constant
is the first higher-order correction to this cubic commutator.
\end{calculation}
\begin{proof}
In \eqref{eq:vol2-moyal}, the odd orders are $n=1$ and $n=3$.  The first contributes $9\hbarq z_1^2z_2^2$.  The third contributes $2(\hbarq/2)^3(3!)=3\hbarq^3/2$.
\end{proof}

\section{Weyl generators}
With the convention $\{q,p\}=1$, the symbols satisfy
$[q,p]_\star=\hbarq$.
Their Weyl-ordered differential operators are
\[
 \widehat q=q,\qquad \widehat p=-\hbarq\partial_q,
 \qquad [\widehat q,\widehat p]=\hbarq.
\]
The abstract Weyl algebra is defined by this commutator relation.
Weyl ordering transports its multiplication to the Moyal product.
The differential-operator representation realizes these relations over
$\C[\hbarq]$ and its formal completion.
The finite quotient argument above uses the abstract $R_N$-algebra,
so it does not require a faithful differential-operator representation at nilpotent $\hbarq$.

\chapter{Continuous Hamiltonians and symplectic descent}
\section{Residue duality and the direction of truncation}
The finite jet algebras are quotients of
$\h=\mathfrak m^2\subset\C\llbracket z_1,z_2\rrbracket$.
Their duals embed into the continuous dual of $\h$.
Continuous duals use the adic topology and the discrete coefficient field.
To see the different directions explicitly, put
\[
 \rho_{ab}=z_1^{-a-1}z_2^{-b-1}\,\dd z_1\dd z_2,
 \qquad a,b\ge0.
\]
The residue pairing is
\[
 \operatorname{Res}_0(z_1^iz_2^j\rho_{ab})
 =\delta_{ia}\delta_{jb}.
\]
Every continuous functional on the power-series ring depends on
finitely many Taylor coefficients.  Hence its continuous dual is
$\bigoplus_{a,b\ge0}\C\rho_{ab}$, and
\[
 \h^\vee_{\mathrm{cont}}
 =\bigoplus_{a+b\ge2}\C\rho_{ab}
 =\varinjlim_d(\h_{\le d})^\vee.
\]
The two degree-one residue modes pair with the translations in the
affine Hamiltonian algebra.

\begin{proposition}[Coadjoint residue formula]
For $f=z_1^iz_2^j\in\h$,
\[
 \ad_f^*\rho_{ab}
 =-\bigl(i(b+1)-j(a+1)\bigr)
       \rho_{a-i+1,b-j+1}.
\]
A mode with a negative index or total index less than two is interpreted
as zero on $\h$.  For a formal $f$, only finitely many of its monomials
contribute to the action on a fixed residue mode.
\end{proposition}
\begin{proof}
By definition,
$\langle\ad_f^*\rho_{ab},g\rangle
=-\langle\rho_{ab},\{f,g\}\rangle$.
For $g=z_1^uz_2^v$, the bracket is
$(iv-ju)z_1^{i+u-1}z_2^{j+v-1}$.
The residue can be nonzero only for
$u=a-i+1$, $v=b-j+1$; substitution gives the stated coefficient.
Nonnegative output indices bound $i\le a+1$ and $j\le b+1$,
which proves finiteness for a power series.
\end{proof}

The continuous semidirect product
\[
 \g=\h\ltimes\h^\vee_{\mathrm{cont}}[-1]
\]
has its invariant evaluation pairing with values in $\C[-1]$.
Give $\h$ its inverse-limit
topology and the residue dual the discrete direct-sum topology.
The coadjoint action is continuous by the finite dependence just proved.
The finite pairs $\g_d$ remain valid cotangent dg Lie algebras, but
they are not the images of a componentwise Lie quotient of $\g$.

\begin{proposition}[Failure of componentwise cotangent truncation]
Let $d\ge2$.  Project $\h$ to degrees at most $d$, and project its
residue dual by retaining only modes with $a+b\le d$.
Their direct sum is not a Lie morphism $\g\to\g_d$.
\end{proposition}
\begin{proof}
Both $z_1^{d+1}$ and $\rho_{d,1}$ project to zero.  Their coadjoint
bracket is
\[
 \ad_{z_1^{d+1}}^*\rho_{d,1}
 =-2(d+1)\rho_{0,2},
\]
which survives the projection.  Thus the kernel is not a Lie ideal.
The natural dual map of the Hamiltonian quotient is instead the
inclusion $(\h_{\le d})^\vee\hookrightarrow\h^\vee_{\mathrm{cont}}$.
\end{proof}

\section{Jet completion and comparison data}
The Hamiltonian Lie algebra and its completed Rees enveloping algebra
do have the quotient systems
\[
 \h=\varprojlim_d\h_{\le d},\qquad
 \widehat U_\hbar(\h)=\varprojlim_d U_\hbar(\h_{\le d}).
\]
This is the jet-quotient topology on the enveloping algebra, together
with its $\hbarq$ topology.  It does not complete arbitrary PBW powers
of a fixed finite jet: for example, a bare infinite sum of powers of
one cubic generator is not an element of $U_\hbar(\h_{\le3})$.
A Lie quotient does not by itself define a covariant map between
Chevalley cochains with adjoint coefficients.
Nor does an algebra quotient define a map between all Hochschild
cochains with values in the algebra itself.
A cochain must descend through the input quotient and the output quotient.
Consequently the coefficient limit of
Corollary~\ref{cor:vol2-formal-relative-rees} is distinct from this jet limit.

A \emph{compatible jet comparison} consists of the following data.
Choose complete $R$-linear cochain models $P_d,H_d$ with transition maps,
levelwise comparisons $P_d\simeq\mathcal P_d$ and
$H_d\simeq C^\bullet_{\HH,\mathrm{cont}}(A_d/R,A_d)$,
and compatible quasi-isomorphisms $P_d\to H_d$.
Use degreewise surjective transitions and ordinary inverse limits,
or use derived inverse limits throughout the following datum.
Choose a continuous Hamiltonian polyvector model $\mathcal P_\infty$,
including its underlying topological space and the differential
$[\hbarq\pi,-]$ for its linear Poisson tensor.
Specify identifications
\[
 \varprojlim_dP_d\simeq\mathcal P_\infty,
 \qquad
 \varprojlim_dH_d\simeq
 C^\bullet_{\HH,\mathrm{cont}}
 (\widehat U_\hbar(\h)/R,\widehat U_\hbar(\h)),
\]
The topology on the right is the stated jet-quotient topology together
with the $\hbarq$ completion.
The identifications must use this same topology.
For higher operations, also supply compatible homotopy-$\Etwo$ and
Swiss-cheese refinements whose operations converge in that topology.

\begin{theorem}[Continuous comparison from compatible jets]\label{thm:vol2-compatible-jets}
A compatible jet comparison induces a continuous quasi-isomorphism
\[
 \Phi_\infty:\mathcal P_\infty
 \xrightarrow{\sim}
 C^\bullet_{\HH,\mathrm{cont}}
 (\widehat U_\hbar(\h)/R,\widehat U_\hbar(\h)).
\]
Compatible homotopy-$\Etwo$ and Swiss-cheese refinements give the
corresponding complete open--closed operations.
\end{theorem}
\begin{proof}
For surjective transitions, the cones of the levelwise maps are
acyclic inverse systems with surjective maps in each degree.
The primitive-lifting argument of
Corollary~\ref{cor:vol2-formal-relative-rees} proves exactness of their limit.
For a derived inverse limit, use its invariance under levelwise equivalences.
Apply the supplied identifications of the two limits.
Compatibility and convergence of the operadic maps give the final assertion.
\end{proof}

The finite PBW comparison constructs the relative centre at each fixed jet.
A continuous field-theory comparison additionally requires the jet datum
and a map from the specified field observables to $\mathcal P_\infty$.
Universal graph formulas alone do not provide the cochain transition maps.

\begin{corollary}[Continuous open--closed structure with compatible jets]
Under the jet comparison and mixed refinement hypotheses of
Theorem~\ref{thm:vol2-compatible-jets}, the pair
$(\mathcal P_\infty,\widehat U_\hbar(\h))$
has the corresponding complete homotopy Swiss-cheese action.
Its closed color is equivalent to the continuous relative derived centre.
\end{corollary}
\begin{proof}
The supplied finite operations commute with the transition maps
and converge in the specified completion.
Their limit acts on the two limit objects.
The comparison theorem identifies the closed complex with continuous relative Hochschild cochains.
\end{proof}


\section{Holomorphic deformation-quantization algebroids}
The symplectic surface gives a second geometric construction.
Let $(S,\omega)$ be a complex surface with a nowhere-degenerate closed
holomorphic two-form, and let $\pi=\omega^{-1}$.
A deformation-quantization algebroid on $S$ is an $R$-linear stack
locally represented by a sheaf of star-product algebras, with
\[
 f\star g-g\star f=\hbarq\{f,g\}_\pi+O(\hbarq^2).
\]
Local algebra identifications and their inner coherence maps are part
of the stack.  A global object, which would present the algebroid by
one global algebra sheaf, is additional data.

\begin{theorem}[Symplectic quantization algebroid]
A holomorphic symplectic surface admits a deformation-quantization
algebroid $\mathcal W_{S,\hbar}$ whose classical limit is the trivial
algebroid of $\Ocal_S$.  With the classical trivialization in the
symplectic classification, its Fedosov class has the form
\[
 \Theta(\mathcal W_{S,\hbar})
 \in\frac{[\omega]}{\hbarq}
       +H^2(S;\C\llbracket\hbarq\rrbracket).
\]
Changing this class can change the equivalence class of the algebroid.
\end{theorem}
\begin{proof}
Apply the symplectic algebroid existence and classification theorems
of \cite{BresslerRojasDQ}*{Theorems~7.1--7.2 and Section~7.3}
to the complex distribution $T^{0,1}S$.
Its invariant functions are $\Ocal_S$; the local Dolbeault lemma gives
the required exactness along the distribution.  The holomorphic
symplectic tensor is the nondegenerate Poisson tensor in that theorem.
The canonical quantization of the trivial classical gerbe supplies
the asserted algebroid.  The Fedosov class records the difference
from this canonical quantization.
\end{proof}

On a formal Darboux disk the completed local algebra is a formal Weyl algebra.
A Fedosov resolution expresses its differential in the form
\[
 D_F=\nabla+\frac1{\hbarq}[\gamma,-],
\]
where the required fibrewise differential may be included in $\gamma$.
The square is an inner derivation by the curvature, so central
curvature gives $D_F^2=0$ while still contributing to the
quantization class.  Changing that central class
is therefore not the same operation as changing a resolution within
a fixed class.

A global holomorphic symplectic connection cannot be assumed even on
a K3 surface.  Its tangent bundle has second Chern number $24$.
A holomorphic connection would have curvature of type $(2,0)$, so
its second Chern form would have type $(4,0)$ and vanish on a complex
surface, a contradiction.  Local holomorphic constructions, or
Dolbeault resolutions with their descent data, are consequently
essential.  Covariant formality for smooth manifolds
\cite{DolgushevCov} uses a chosen connection; equivariance there
requires an invariant connection.

Holomorphic Lie algebroid formality also gives a sheafwise
$L_\infty$ comparison between polyvectors and polydifferential
operators, together with the compatible comparison of their
differential-form and Hochschild-chain modules
\cite{CalaqueDolgushevHalbout}*{Theorem~5.2 and Corollary~5.3}.
For the tangent Lie algebroid, this supplies holomorphic formality
without a global holomorphic connection.  The additional
$\Etwo$ structure and descent for a specified quantization class
are the data in the following comparison.

\section{Centre descent with a fixed formal Poisson class}
All cochains in this section are continuous and relative to
$R=\C\llbracket\hbarq\rrbracket$.
Let $\{U_i\}$ be a Darboux cover and choose a formal Poisson tensor
\[
 \Pi_i=\hbarq\pi_i+\hbarq^2\pi_{i,2}+\cdots,
 \qquad [\Pi_i,\Pi_i]_{\mathrm{Sch}}=0,
\]
whose quantization represents the chosen local algebroid.
Write $\PV_{U_i}$ for holomorphic polyvector fields, placing
$\Lambda^kT_{U_i}$ in degree $k$.
Fix homotopy-$\Etwo$ structures and local equivalences
\[
 \widetilde\Phi_i:
 (\PV_{U_i}\llbracket\hbarq\rrbracket,[\Pi_i,-])
 \simeq C^\bullet_{\HH,\mathrm{cont}}
       (\mathcal W_{U_i,\hbar}/R,\mathcal W_{U_i,\hbar}).
\]
Require compatible equivalences on overlaps and their coherent
homotopies on all higher intersections, both for the source complexes
and for these maps.  Denote the resulting sheaf of source algebras
by $\mathcal P_{S,\Pi}$ and the algebroid's relative Hochschild
cochain sheaf by $\mathcal C_{S,\hbar}$.

\begin{theorem}[Global centre comparison with coherent descent]
The specified data induce a sheafwise homotopy-$\Etwo$ equivalence
\[
 \mathcal P_{S,\Pi}\simeq\mathcal C_{S,\hbar}
\]
and an equivalence of their derived global sections.
If the chosen continuous Morita theory identifies
$Z_{\mathrm{Mor},R}^{\mathrm{der}}(\Perf(\mathcal W_{S,\hbar}))$
with $\R\Gamma(S,\mathcal C_{S,\hbar})$, this gives
\[
 Z_{\mathrm{Mor},R}^{\mathrm{der}}
       (\Perf(\mathcal W_{S,\hbar}))
 \simeq\R\Gamma(S,\mathcal P_{S,\Pi}).
\]
\end{theorem}
\begin{proof}
The overlap maps and their higher homotopies define a map of descent
objects.  Its restriction to each $U_i$ is an equivalence, so it is
an equivalence of sheaves.  Derived global sections preserve this
equivalence.  Applying the stated Morita identification gives the
last formula.
\end{proof}

When the chosen formal Poisson class is represented globally by
$\Pi=\hbarq\pi$ with the indicated untwisted descent, the source is
$(\PV_S\llbracket\hbarq\rrbracket,[\hbarq\pi,-])$.
An arbitrary quantization period does not specify this representative.
Likewise, equivariance under holomorphic symplectomorphisms requires
equivariant comparison and descent data; local formality alone does
not supply them.

\section{Topological factorization with holomorphic coefficients}
Give $\R^2$ its standard framing.  A complete homotopy-$\Etwo$
algebra internal to sheaves on $S$ has a locally constant
factorization envelope on $\R^2$ valued in those sheaves, using the
chosen completed tensor product.  Applied to the preceding equivalence,
this gives
\[
 \operatorname{Fact}_{\Etwo}(\mathcal P_{S,\Pi})
 \simeq\operatorname{Fact}_{\Etwo}(\mathcal C_{S,\hbar}).
\]
The assertion follows by applying the factorization envelope to an
equivalence of its coefficient algebras.
It retains the relative centre action on the corresponding local
perfect-module categories.

The holomorphic variable here is a sheaf variable: restriction to
smaller open sets is contravariant.  A factorization algebra on
$\R^2\times S$ also requires covariant operations for disjoint
holomorphic open sets and the relevant descent condition.
Such operations require a geometric construction in the coefficient
direction.  Once supplied compatibly on both sides, they promote the
displayed equivalence to that geometric factorization comparison.
They are not part of the definition of an $\Etwo$-algebra sheaf.

\section{Classical mixed fields and the quantization question}
There is a local classical BF construction with the desired mixed
directions.  On a holomorphic symplectic chart $U\subset S$, let
\[
 \mathcal D=\Omega^\bullet(\R^2)\widehat\tensor
                \Omega^{0,\bullet}(U),
 \qquad D=\dd_{\R^2}+\bar\partial.
\]
Here the tensor is the completed projective tensor of smooth-form
spaces.  The differential acts by the graded tensor rule.
For a finite-dimensional Lie algebra $L$, take fields
\[
 A\in\mathcal D\tensor L[1],\qquad
 B\in\mathcal D\tensor L^\vee[2].
\]
The two de Rham and two Dolbeault degrees make the integration degree
four.  The shifts above therefore give the BV pairing degree $-1$.
With compact support, the degree-zero action is
\[
 S_{\mathrm{BF}}(A,B)
 =\int_{\R^2\times U}\omega\wedge
     \left\langle B,DA+\frac12[A,A]\right\rangle.
\]
Its classical master equation follows from $D^2=0$, the derivation
property of $D$, Jacobi, and integration by parts.
For $L=\h_{\le d}$ this is a finite coefficient model.
A formal Hamiltonian field theory must additionally specify its
continuous coefficient complex and the gluing of formal coordinate
charts.

Renormalized quantum observables require compatible gauge fixing and
counterterms satisfying the field-theoretic quantum master equation.
The finite identity $\Delta S_L=0$ controls its stated finite model;
it does not establish the local field-theoretic quantum master equation.
Even after quantization, identifying the observable algebra with the
relative centre requires a map from those observables to the selected
Hamiltonian or symplectic polyvector complex, with the factorization
operations and descent respected.

\chapter{Derived matrix equations and necklace actions}
\section{The derived commuting fibre}
For a positive integer $N$, let $X=(X_{ij})$ and $Y=(Y_{ij})$ be
two generic $N\times N$ matrices.  Define the commutative dg algebra
\[
 \mathcal K_N=
 \C[X_{ij},Y_{ij}]\tensor\Lambda(\psi_{ij}),
 \qquad |X_{ij}|=|Y_{ij}|=0,\quad |\psi_{ij}|=-1,
\]
with differential
\[
 QX=QY=0,\qquad Q\psi=[X,Y].
\]
Here $\psi$ is a matrix with anticommuting entries, and products of
such matrices use the ordinary matrix rule with graded coefficients.
The algebra $\mathcal K_N$ is the function algebra of the derived
zero fibre of the moment map
\[
 \mu:\mathfrak{gl}_N^2\longrightarrow\mathfrak{gl}_N,
 \qquad (X,Y)\longmapsto[X,Y].
\]
Indeed, the Koszul resolution of the origin in the target, pulled back
along $\mu$, computes the derived tensor product defining this fibre.
Its zeroth cohomology is the coordinate algebra of the ordinary
commuting scheme.
Its negative cohomology retains the failure of the commuting equations
to form a regular sequence.

For example, $\Tr\psi$ is a closed element of degree $-1$ because
$\Tr[X,Y]=0$.  It is not exact: setting $X=Y=0$ annihilates every
boundary but leaves the nonzero exterior element $\Tr\psi$.
At rank one this becomes
\[
 \mathcal K_1=\C[x,y]\tensor\Lambda(\psi),\qquad Q=0.
\]
Thus even the simplest derived commuting fibre has an odd direction
in addition to its ordinary affine plane.

Simultaneous conjugation acts on $\mathcal K_N$.
Two quotients of this action will be distinguished:
\[
 [\operatorname{Spec}\mathcal K_N/GL_N],
 \qquad
 \mathcal B_N=C^*_{\CE}(\mathfrak{gl}_N,\mathcal K_N).
\]
The first is the derived group quotient; the second computes
Chevalley cochains for the infinitesimal action.
They need not have the same function complex.
At rank one conjugation is trivial, and rational $\mathbb G_m$
invariants are exact, so the group quotient has function complex
$\mathcal K_1$.  In contrast,
\[
 \mathcal B_1=\mathcal K_1\tensor\Lambda(c),\qquad |c|=1,
\]
with zero differential.  This additional Chevalley ghost records
Lie algebra cohomology.

\section{Cyclic words and the moment map}
Let $F=\C\angles{x,y}$ and let $F_{\mathrm{cyc}}=F/[F,F]$,
where $[F,F]$ is the linear span of associative commutators.
For a word $a_1\cdots a_r$, its cyclic derivative in $x$ is
\[
 \partial_x^{\mathrm{cyc}}[a_1\cdots a_r]
 =\sum_{a_i=x}a_{i+1}\cdots a_ra_1\cdots a_{i-1}.
\]
Define $\partial_y^{\mathrm{cyc}}$ similarly.
These derivatives are well defined on cyclic words and take values
in the free algebra.  Put
\[
 \delta_f(x)=-\partial_y^{\mathrm{cyc}}f,
 \qquad \delta_f(y)=\partial_x^{\mathrm{cyc}}f,
\]
and extend by the derivation rule.

\begin{theorem}[Necklace Lie algebra and action]
The bracket
\[
 \{f,g\}_{\mathrm{neck}}
 =\left[
   \partial_x^{\mathrm{cyc}}f\,\partial_y^{\mathrm{cyc}}g
  -\partial_y^{\mathrm{cyc}}f\,\partial_x^{\mathrm{cyc}}g
  \right]_{\mathrm{cyc}}
\]
makes $F_{\mathrm{cyc}}$ a Lie algebra, with the constants central.
The map $f\mapsto\delta_f$ is a Lie morphism, factors through
$F_{\mathrm{cyc}}/\C$, and satisfies
\[
 \delta_f([x,y])=0.
\]
\end{theorem}
\begin{proof}
Consider the bilinear map $F\tensor F\to F\tensor F$ determined by
the generator values
$\{\!\{x,y\}\!\}=1\tensor1$,
$\{\!\{y,x\}\!\}=-1\tensor1$, with the other generator brackets
zero, the outer Leibniz rule in the second argument, and skew-symmetry
with interchange of the two output factors.
This is a double bracket.  Its double Jacobiator is the cyclic sum
of nested brackets, with the outer bracket acting in the first output
factor.  It vanishes on generator triples, since the inner brackets
are constant.  The derivation property in each argument then gives
its vanishing on all words.
Its induced bracket on cyclic words is the displayed formula.
The double Jacobi identity gives both the Lie identity and
$[\delta_f,\delta_g]=\delta_{\{f,g\}_{\mathrm{neck}}}$;
see \cite{VanDenBerghDouble}*{Section~2.4 and Theorem~6.3.1}.
The cyclic derivative identity
\[
 [x,\partial_x^{\mathrm{cyc}}f]
 +[y,\partial_y^{\mathrm{cyc}}f]=0
\]
follows by cancellation of successive cyclic cuts of each word.
Substitution into the derivative of $[x,y]$ proves the moment-map
identity.  Constant cyclic words have zero derivatives.
\end{proof}

\begin{proposition}[Action on the derived commuting fibre]
Matrix evaluation of $\delta_f$ on $X,Y$, together with
$\delta_f\psi=0$, defines an action of $F_{\mathrm{cyc}}/\C$
on $\mathcal K_N$ by dg derivations.
It is equivariant for conjugation and hence acts on
$\mathcal B_N$ as well.
\end{proposition}
\begin{proof}
The moment-map identity gives
$\delta_f(Q\psi)=\delta_f[X,Y]=0=Q(\delta_f\psi)$.
The commutator with $Q$ also vanishes on $X,Y$, so it vanishes on
the whole algebra.  Matrix evaluation preserves the free-algebra
derivation and Lie identities.
All expressions are equivariant polynomials in the matrices;
extending the derivation by zero on Chevalley ghosts therefore gives
the action on $\mathcal B_N$.
\end{proof}

The ordinary matrix space carries the symplectic form
$\Tr(\dd X\wedge\dd Y)$.
For $f\in F_{\mathrm{cyc}}$ define $J_N(f)=\Tr f(X,Y)$.
Constants are retained here: $J_N(1)=N$.

\begin{theorem}[Trace moment map]
With this symplectic convention,
\[
 \{J_N(f),J_N(g)\}_{\mathfrak{gl}_N^2}
 =J_N(\{f,g\}_{\mathrm{neck}}).
\]
\end{theorem}
\begin{proof}
Cyclic differentiation gives
\[
 \dd J_N(f)=\Tr\bigl(
       (\partial_x^{\mathrm{cyc}}f)(X,Y)\dd X
      +(\partial_y^{\mathrm{cyc}}f)(X,Y)\dd Y\bigr).
\]
Contraction with the inverse trace symplectic form gives the necklace
formula.  Equivalently, the matrix-entry brackets are
$\{X_{ij},Y_{kl}\}=\delta_{il}\delta_{jk}$,
as in \cite{VanDenBerghDouble}*{Propositions~7.5.1--7.5.2}.
\end{proof}

\section{Formal completion and the enveloping action}
Let $F_{\mathrm{cyc},r}$ be the cyclic words of length $r$ and put
\[
 \mathfrak n=\prod_{r\ge2}F_{\mathrm{cyc},r}.
\]
Its bracket sends lengths $r,s$ to length $r+s-2$.
Complete $\mathcal K_N$ by the weights
\[
 \operatorname{wt}X_{ij}=\operatorname{wt}Y_{ij}=1,
 \qquad \operatorname{wt}\psi_{ij}=2.
\]
The differential preserves weight.
For a Hamiltonian of length $r$, its derivation raises weight by
$r-2$ on the even generators and annihilates $\psi$.
The sum of these derivations therefore converges on the weight
completion $\widehat{\mathcal K}_N$.
Chevalley ghosts have weight zero, giving the corresponding
completion $\widehat{\mathcal B}_N$.

Let $\mathfrak n_{\le d}=\bigoplus_{2\le r\le d}F_{\mathrm{cyc},r}$
with the projected bracket, and define
$\widehat U_\hbar(\mathfrak n)=\varprojlim_d
U_\hbar(\mathfrak n_{\le d})$.
The continuous Lie action gives a Rees enveloping action
\[
 \widehat U_\hbar(\mathfrak n)
 \longrightarrow
 \operatorname{End}_{\mathrm{cont}}
       (\widehat{\mathcal B}_N)\llbracket\hbarq\rrbracket,
 \qquad f\longmapsto\hbarq\delta_f.
\]
The enveloping completion uses these length quotients and the
coefficient parameter.  A Hamiltonian of length greater than $d$
acts trivially on the quotient by matrix weights at least $d$.
Thus the finite actions are compatible on every bounded matrix
weight quotient, and their inverse limit defines the displayed map.
The Rees relation follows directly from
\[
 [\hbarq\delta_f,\hbarq\delta_g]
 =\hbarq^2\delta_{\{f,g\}_{\mathrm{neck}}}.
\]
At any fixed output weight only finitely many Hamiltonian lengths
contribute.  The topology on endomorphisms is the inverse-limit
topology of these weight quotients.  No completion by arbitrary PBW
powers in a fixed finite length quotient is imposed.

\section{The commutative Hamiltonian lift}
Abelianization of words gives a continuous surjective Lie map
\[
 \mathfrak n\longrightarrow\h.
\]
Here $x$ maps to $z_1$ and $y$ to $z_2$.
A linear section is cyclic symmetrization: average the words with a
prescribed number of $x$'s and $y$'s and pass to their cyclic class.
Denote it by $\sigma$.  It need not be a Lie section.

\begin{proposition}[The first symmetrization defect]
For $f=x^3$ and $g=y^3$,
\[
\begin{split}
 D&=\{\sigma f,\sigma g\}_{\mathrm{neck}}
                      -\sigma\{f,g\}\\
  &=3\bigl([xxyy]-[xyxy]\bigr)
   =-\frac32\bigl[[x,y]^2\bigr]_{\mathrm{cyc}}\ne0.
\end{split}
\]
On the derived commuting fibre, the corresponding derivation defect
is null-homotopic.  A degree-$-1$ derivation giving the homotopy is
\[
 H(X)=-3[X,\psi],\qquad
 H(Y)=-3[Y,\psi],\qquad
 H(\psi)=3\psi^2.
\]
It satisfies $QH+HQ=\delta_D$ and extends equivariantly to
$\mathcal B_N$.
\end{proposition}
\begin{proof}
There are six words with two copies of each generator.  Four have cyclic
class $[xxyy]$ and two have cyclic class $[xyxy]$, so
\[
 \sigma(x^2y^2)=\frac23[xxyy]+\frac13[xyxy].
\]
The necklace bracket of $[x^3]$ and $[y^3]$ is $9[xxyy]$,
whereas the commutative bracket is $9x^2y^2$.
Their difference is the first displayed expression for $D$.
Expansion of $[x,y]^2$ gives
$2[xyxy]-2[xxyy]$ after cyclic reduction.
The two necklaces are distinct basis classes, proving nonvanishing.
The derivation defect is also nonzero at rank two: for
$X=\operatorname{diag}(1,0)$ and $Y=E_{12}$ one has
$[X,[X,Y]]=E_{12}$, hence $\delta_D(X)=-3E_{12}$.

Put $\mu=[X,Y]$.  Cyclic differentiation of $D$ gives
\[
 \delta_D(X)=-3[X,\mu],\qquad
 \delta_D(Y)=-3[Y,\mu],\qquad \delta_D(\psi)=0.
\]
The first two identities for $QH+HQ$ follow from $Q\psi=\mu$.
On the odd generator, the graded derivation rule gives
\[
 QH(\psi)=3(\mu\psi-\psi\mu),\qquad
 HQ(\psi)=-3(\mu\psi-\psi\mu),
\]
so these terms cancel.  Equality on generators gives equality on
the whole free graded algebra and hence after matrix evaluation.
The formulas are conjugation equivariant, so extending $H$ by zero
on the Chevalley ghosts preserves the same homotopy identity.
\end{proof}

This homotopy resolves one binary defect.  An action of all
commutative formal Hamiltonians on the derived matrix object requires
higher homotopies satisfying the $\Linf$ identities and converging
in the specified topology.  Constructing such a lift, or a different
strict Lie lift, remains the comparison problem.
The necklace action is already a complete positive construction;
the failed strictness of $\sigma$ is not an obstruction to every
possible Hamiltonian action.

\section{The first two ordinary trace models}
Passing to ordinary rank-one functions gives $\C[x,y]$ and
$J_1(f)=f(x,y)$ after abelianization.
The necklace bracket becomes the usual Poisson bracket.
Pointed Hamiltonians act by vector fields fixing the origin; linear
Hamiltonians generate translations.  The odd coordinate of
$\mathcal K_1$ and the Chevalley ghost of $\mathcal B_1$ belong to
their larger derived models.

For two-by-two matrices, Cayley--Hamilton gives
\[
 X^2-(\Tr X)X+(\det X)I=0
\]
and the analogous identity for $Y$.
These reduce a consecutive power $X^2$ or $Y^2$ in a word after
adjoining the indicated trace and determinant coefficients.
Commuting complex matrices admit a common invariant flag and hence
can be simultaneously triangularized.  Their traces are therefore
computed from the diagonal pairs $(x_i,y_i)$, even when the matrices
are not simultaneously diagonalizable.
Write $p_{ab}=\sum_{i=1}^2x_i^ay_i^b$.
Then
\[
 p_{20}=p_{10}^2-2\det X,\qquad
 p_{02}=p_{01}^2-2\det Y,
\]
\[
 p_{11}^2-p_{20}p_{02}
 =-(x_1y_2-x_2y_1)^2.
\]
The first two identities are Cayley--Hamilton trace identities; the
third follows by expanding the two diagonal sums.
These are identities of ordinary trace functions.  They do not
identify the full derived quotient with its eigenvalue model.

\chapter{Stable traces and the centre comparison}
\section{Graded stabilization}
The eigenvalue model of rank $N$ has coordinate ring
\[
 \mathcal T_N=
 \C[x_1,y_1,\ldots,x_N,y_N]^{S_N}.
\]
Give $x_i,y_i$ weight one.  For $a,b\ge0$ with $a+b>0$, put
\[
 p_{ab}^{(N)}=\sum_{i=1}^Nx_i^ay_i^b.
\]
Adding the zero pair defines restriction maps
$\mathcal T_{N+1}\to\mathcal T_N$ preserving this weight.
Define the graded stable algebra and its weight completion by
\[
 \mathcal T_{\mathrm{gr}}
 =\bigoplus_{w\ge0}\varprojlim_N(\mathcal T_N)_w,
 \qquad
 \widehat{\mathcal T}_{\mathrm{wt}}
 =\prod_{w\ge0}(\mathcal T_{\mathrm{gr}})_w.
\]

\begin{theorem}[Stable multisymmetric algebra]
For each weight $w$, restriction stabilizes once $N\ge w$.
There is a graded algebra isomorphism
\[
 \mathcal T_{\mathrm{gr}}
 \cong\C[p_{ab}\mid a,b\ge0,\ a+b>0],
 \qquad \operatorname{wt}p_{ab}=a+b.
\]
Its weight completion is $\widehat{\mathcal T}_{\mathrm{wt}}$.
\end{theorem}
\begin{proof}
A monomial orbit sum is indexed by a multiset of nonzero exponent
pairs $\alpha_1,\ldots,\alpha_l$, placed in distinct slots among
$1,\ldots,N$.  At total weight $w$ one has $l\le w$.
For $N\ge w$ all such multisets occur, and zero-pair restriction
identifies their orbit-sum bases.  This proves stabilization.

Expand $p_{\alpha_1}\cdots p_{\alpha_l}$ by assigning a slot to each
factor.  The assignments with distinct slots give the corresponding
orbit sum with coefficient $\prod_\alpha m_\alpha!$, where
$m_\alpha$ is the multiplicity of $\alpha$ in the multiset.
All other assignments merge two or more exponent pairs and have
fewer occupied slots.  Induction on $l$ therefore gives a triangular
change of basis between the power-sum monomials and the orbit sums.
The diagonal coefficients are nonzero over $\C$.
In each weight the two bases are consequently equivalent, proving
generation and algebraic independence of the stable power sums.
Taking the product over weights gives the stated completion.
\end{proof}

This definition distinguishes graded stabilization from an inverse
limit of ungraded polynomial rings.  For example,
$\sum_{r\ge0}p_{10}^r$ is a valid weight-completed series but does
not restrict to a polynomial at rank one.  It is therefore not an
element of $\varprojlim_N\mathcal T_N$ with uncompleted terms.
Finite-rank relations and the formal topology must both be specified
before taking a stable limit.

\section{The rank parameter and the pointed Poisson algebra}
On $(\C^2)^N$ use $\{x_i,y_j\}=\delta_{ij}$.
The invariant bracket is
\[
 \{p_{ab}^{(N)},p_{cd}^{(N)}\}
 =(ad-bc)p_{a+c-1,b+d-1}^{(N)},
 \qquad p_{00}^{(N)}=N.
\]
It follows by differentiating the two sums; terms with distinct
indices have zero bracket.  A negative output index occurs only
with zero coefficient.
In particular,
$\{p_{10}^{(N)},p_{01}^{(N)}\}=N$.
Adding the zero pair is therefore not Poisson on the full trace algebra.

There are two compatible ways to retain this formula.
First, introduce a central symbol $\nu$ and define
\[
 \widetilde{\mathcal T}
 =\C[\nu,p_{ab}\mid a+b>0],\qquad p_{00}=\nu,
\]
with the displayed bracket on generators and $\{\nu,-\}=0$.
This is the symmetric Poisson algebra of the Lie algebra of polynomial
Hamiltonians including its constant.  Jacobi follows from the
polynomial Poisson bracket.  Specialization $\nu=N$ and
$p_{ab}\mapsto p_{ab}^{(N)}$ is a Poisson map for every rank.

Second, use the pointed subalgebra generated by $p_{ab}$ with
$a+b\ge2$.
The bracket of two such generators again has degree at least two,
so it never produces the rank parameter.  Zero-pair restriction is
Poisson on this subalgebra.  Its weight completion is also Poisson:
the bracket lowers weight by two, and only finitely many homogeneous
input weights can contribute to a fixed output weight.
This gives the stable pointed Hamiltonian trace algebra.

\section{The two-brane cone}
Separate center-of-mass and relative coordinates.  Put
\[
 u=x_1-x_2,
 \qquad
 v=y_1-y_2,
 \qquad
 E=\frac12u^2,
 \quad
 F=-\frac12v^2,
 \quad
 H=-uv.
\]
Then
\[
 H^2+4EF=0.
\]

\begin{proposition}[The $A_1$ Kirillov--Kostant cone]
With $\{u,v\}=2$,
\[
 \{H,E\}=4E,
 \qquad
 \{H,F\}=-4F,
 \qquad
 \{E,F\}=2H.
\]
The generators $E/2,F/2,H/2$ satisfy the standard
$\mathfrak{sl}_2$ Lie--Poisson brackets.  Their quadratic relation
defines its nilpotent cone.
\end{proposition}
\begin{proof}
The three identities follow by direct differentiation.  The quadratic equation is the determinant-zero equation for the traceless matrix
\[
 \begin{pmatrix}H&2E\\2F&-H\end{pmatrix}.
\]
\end{proof}

\section{Reduced cyclic traces and stable Lie homology}
A trace of the unit records the representation dimension.
For $B=\C$, one has $\operatorname{Tr}_N(1)=N$, so ordinary traces
cannot define a rank-compatible map on all of $HC_0(B)$.
A scalar-reduced trace removes this dependence after choosing an augmentation.

Let $B$ be a unital associative $\C$-algebra.
The affine representation scheme $\Rep_NB$ represents unital algebra maps
$B\to\operatorname{Mat}_N(D)$ for commutative $\C$-algebras $D$.
Put
\[
 \mathcal R_N=\Ocal(\Rep_NB)^{GL_N},\qquad
 \mathfrak c_B=HC_0(B)=B/[B,B],
\]
where $[B,B]$ denotes the linear span of commutators.
The ordinary trace is the linear map
\[
 \operatorname{Tr}_N:\mathfrak c_B\longrightarrow\mathcal R_N,
 \qquad
 \operatorname{Tr}_N([b])(\rho)=\operatorname{Tr}(\rho(b)).
\]
Cyclicity makes it independent of the representative, and conjugation
invariance puts its image in $\mathcal R_N$.

For Lie and Poisson compatibility, specify a Lie bracket on $\mathfrak c_B$,
a Poisson bracket on each $\mathcal R_N$, and the identities
\begin{equation}\label{eq:vol2-trace-bracket-datum}
 \{\operatorname{Tr}_N\beta,\operatorname{Tr}_N\gamma\}_N
 =\operatorname{Tr}_N[\beta,\gamma]_{\mathrm{cyc}}
 \qquad(\beta,\gamma\in\mathfrak c_B).
\end{equation}
A Hamiltonian comparison also specifies a Lie map
$j:\h\to\mathfrak c_B$ for this bracket.
These are additional data on $B$.
The preceding necklace bracket and trace moment map supply the bracket
and trace data for $B=\C\angles{x,y}$ with its canonical symplectic bracket.
The map $j$ remains an additional comparison.
Here $j$ is an algebraic Lie map.  A continuous formal comparison
also requires a topology on the cyclic recipient and continuity of
$j$ and the trace maps.
An arbitrary unital algebra does not carry these structures by definition.

For rank stabilization, choose an augmentation $\epsilon:B\to\C$.
It defines the rank-stabilization morphism
\[
 s_N:\Rep_NB\longrightarrow\Rep_{N+1}B,
 \qquad \rho\longmapsto\rho\oplus\epsilon.
\]
Restriction gives $s_N^*:\mathcal R_{N+1}\to\mathcal R_N$.
Define the reduced cyclic space and trace by
\[
 \overline{\mathfrak c}_B=B/(\C1+[B,B]),
 \qquad
 t_N(\overline b)(\rho)
 =\operatorname{Tr}(\rho(b))-N\epsilon(b).
\]
Let $\mathfrak m_N\subset\mathcal R_N$ be the ideal of functions
vanishing at the representation $b\mapsto\epsilon(b)I_N$.
The completed invariant rings and stable target are
\[
 \widehat{\mathcal R}_N
 =\varprojlim_{r\ge1}\mathcal R_N/\mathfrak m_N^r,
 \qquad
 \mathcal R_\infty^{\wedge}
 =\varprojlim_N\widehat{\mathcal R}_N.
\]
Use the inverse-limit topologies.
Write $J$ for the augmentation ideal in $\Sym_\C\overline{\mathfrak c}_B$
and put
\[
 \widehat\Sym\overline{\mathfrak c}_B
 =\varprojlim_{r\ge1}
   \Sym_\C\overline{\mathfrak c}_B/J^r.
\]

The stable matrix Lie algebra is
$\mathfrak{gl}_\infty B=\varinjlim_N\mathfrak{gl}_N B$
under block-zero Lie embeddings.
For its homology, use trivial coefficients $\C$.
Let $V_B=\bigoplus_{n\ge1}HC_{n-1}(B)$, with $HC_{n-1}(B)$ placed
in homological degree $n$.
The notation $\Sym_{\mathrm{gr}}V_B$ denotes the free graded-commutative
algebra, with the Koszul sign rule.

\begin{theorem}[Cyclic traces and stable Lie homology]
Assume the compatible jet comparison and its homotopy Gerstenhaber refinement.
Write $\mathcal A=\widehat U_\hbar(\h)$.
For the augmented algebra and completion above, there are maps
\[
 \mathcal P_\infty
 \xrightarrow{\ \Phi_\infty\ }
 C^*_{\HH,\mathrm{cont}}(\mathcal A/R,\mathcal A),
 \qquad
 \operatorname{Tr}_N:\mathfrak c_B\longrightarrow\mathcal R_N,
\]
\[
 \widehat\Sym\overline{\mathfrak c}_B
 \longrightarrow\mathcal R_\infty^{\wedge},
 \qquad
 \Sym_{\mathrm{gr}}V_B
 \xrightarrow{\sim}
 H_*^{\CE}(\mathfrak{gl}_\infty B;\C).
\]
The first preserves the selected homotopy Gerstenhaber structure.
The second is linear and, under the bracket data
\eqref{eq:vol2-trace-bracket-datum}, is a Lie map.
Under those data, its composite with $j$ preserves the Hamiltonian bracket.
The third is a continuous homomorphism of commutative algebras,
and the fourth is an isomorphism of graded Hopf algebras.
These maps are natural for morphisms preserving the stated data.
\end{theorem}
\begin{proof}
The first arrow is the chosen comparison of
Theorem~\ref{thm:vol2-compatible-jets}.
Its operadic refinement gives the asserted Gerstenhaber compatibility.
The trace formula above constructs the second arrow.
Equation~\eqref{eq:vol2-trace-bracket-datum} supplies its Lie compatibility,
and the Lie property of $j$ gives the composite assertion.

The augmentation annihilates commutators.
Both terms of $t_N$ vanish on a commutator, and their values on a scalar
multiple of the unit cancel.
Thus $t_N$ is well defined on $\overline{\mathfrak c}_B$.
At the trivial representation its value is zero, so its image lies in $\mathfrak m_N$.
Moreover,
\[
\begin{split}
 (s_N^*t_{N+1})(\overline b)(\rho)
 &=\operatorname{Tr}(\rho(b))+\epsilon(b)-(N+1)\epsilon(b)\\
 &=t_N(\overline b)(\rho).
\end{split}
\]
The point $N\epsilon$ maps to $(N+1)\epsilon$, so
$s_N^*(\mathfrak m_{N+1}^r)\subset\mathfrak m_N^r$.
Hence restriction defines the transition maps between the completed rings.
The symmetric-algebra extension of $t_N$ sends $J^r$ into $\mathfrak m_N^r$.
It therefore induces a ring homomorphism on each corresponding quotient.
Taking the inverse limit in $r$ gives a continuous map to
$\widehat{\mathcal R}_N$.
The displayed stabilization identity makes these maps compatible in $N$.
Their inverse limit is the third arrow.

The diagonal Lie map
$\mathfrak{gl}_\infty B\to
 \mathfrak{gl}_\infty B\oplus\mathfrak{gl}_\infty B$
induces the homology coproduct.
Block sum in the opposite direction induces the homology product,
using the K\"unneth isomorphism over $\C$.
The Loday--Quillen--Tsygan theorem
\cite{LodayQuillenMatrices}*{Theorem~6.2} identifies primitive homology
in degree $n$ with $HC_{n-1}(B)$.
The stable homology is a connected graded-commutative and cocommutative
Hopf algebra over $\C$.  The characteristic-zero structure theorem
therefore identifies it with the free graded-commutative Hopf algebra
on its primitives, giving the fourth arrow.
The trace formulas, the augmentation-preserving rank maps, and stable
Lie homology are natural.
Naturality of the first arrow requires preservation of the selected jet comparison.
\end{proof}

Scalar reduction need not preserve the supplied Poisson bracket.
For example, for $B=\C[x,y]$ with $\{x,y\}=1$ and
$\epsilon(x)=\epsilon(y)=0$, traces on $N$ points satisfy
\[
 \{\operatorname{Tr}_Nx,\operatorname{Tr}_Ny\}=N,
 \qquad t_N(\overline1)=0.
\]
A reduced Poisson comparison requires the cyclic bracket to descend
through the scalar quotient and to satisfy the reduced trace identity.
For example, centrality of the unit and vanishing of $\epsilon$ on
the selected cyclic brackets give these properties: subtracting
$N\epsilon(b)$ does not change the bracket of two trace functions,
and then $t_N([b,c])=\operatorname{Tr}_N[b,c]$.
A restriction to a Hamiltonian subalgebra may supply these conditions.
Only multiplicativity and continuity are asserted for the third arrow.
For $B=\C$, the reduced cyclic space is zero and that arrow is the identity on $\C$.

\section{The local centre and its modules}
Assume a compatible jet comparison and a complete homotopy-$\Etwo$
refinement as in Theorem~\ref{thm:vol2-compatible-jets}.
Put
\[
 \Acal_{\mathrm{bulk}}^{\mathrm{alg}}
 =\operatorname{Fact}_{\Etwo}(\mathcal P_\infty).
\]
The factorization envelope uses the standard framing of the two topological directions.
The formal holomorphic coefficient remains at the Darboux point.
Let $\Ccal_\partial$ be the category of perfect complete $\mathcal A$-modules
in the chosen continuous $R$-linear Morita theory.
Assume its centre is computed by
$C^\bullet_{\HH,\mathrm{cont}}(\mathcal A/R,\mathcal A)$.

\begin{theorem}[Local derived-centre comparison]
These data give an equivalence of complete homotopy-$\Etwo$ factorization algebras
\[
 \Acal_{\mathrm{bulk}}^{\mathrm{alg}}
 \xrightarrow{\sim}
 \operatorname{Fact}_{\Etwo}\!\left(
 Z^{\mathrm{der}}_{\mathrm{Mor},R}(\Ccal_\partial)\right).
\]
On a topological disk the coefficient map is the selected refinement
of $\Phi_\infty$.  Every perfect complete boundary module has the
induced derived centre action.
\end{theorem}
\begin{proof}
The compatible jet theorem gives the complete homotopy-$\Etwo$ comparison.
The factorization envelope preserves this equivalence.
The assumed continuous Morita model identifies its target with the centre.
A centre element is a derived natural endomorphism of the identity functor,
so evaluation gives the action on each boundary module.
\end{proof}

More generally, a complete module $M$ receives the same action whenever
the continuous Morita theory supplies the completed derived tensor and
$\mathcal A\widehat\tensor_{\mathcal A}^{\mathbf L}M\simeq M$.
Apply $-\widehat\tensor_{\mathcal A}^{\mathbf L}M$ to the derived
bimodule endomorphisms of $\mathcal A$ which define Hochschild cochains.

\begin{corollary}[Global K3 comparison with descent data]
Let $S$ be a K3 surface with holomorphic symplectic form.
Suppose the local comparisons extend, with the coherent descent data
specified above, to a sheafwise equivalence with
$\mathcal C_{S,\hbar}$, and that the stated continuous Morita
identification holds.
Suppose further that the selected cyclic and mixed refinements extend
and satisfy descent in the same completed category.
Then derived global sections give the corresponding global centre,
cyclic, and open--closed comparisons.
\end{corollary}
\begin{proof}
Derived global sections preserve equivalences of sheaves.
Apply them to the supplied sheafwise maps and their coherence homotopies.
\end{proof}


\section{The remaining comparison maps}
At a fixed Hamiltonian jet, PBW antisymmetrization proves the relative
Rees comparison, including both Chevalley terms scaled by $\hbarq$.
Affine formality supplies a comparison with its selected $\Linf$
structure, and Duflo
identifies the degree-zero invariant algebras.  A homotopy-$\Etwo$
or mixed refinement is specified when those operations are used.
The oscillator algebra is a separate associative quotient; its
Casimir relation obstructs descent of the primitive coproduct.

For the Hamiltonian inverse limit, the compatible cochain systems
of Theorem~\ref{thm:vol2-compatible-jets} are the required datum.
With its operadic refinement and continuous Morita model, the
centre acts on every supplied complete boundary module $M$ through
\[
 \mathcal P_\infty
 \xrightarrow{\ \Phi_\infty\ }
 C^\bullet_{\HH,\mathrm{cont}}(\mathcal A/R,\mathcal A)
 \longrightarrow\RHom_{\mathcal A}(M,M).
\]
The final arrow evaluates a derived natural endomorphism at $M$.
For the matrix object $\widehat{\mathcal B}_N$, the necklace action
has been constructed.  A coherent lift of commutative Hamiltonians
would provide the additional action needed to compare it with a
module over the Hamiltonian boundary algebra.
The explicit binary homotopy gives one equation in this problem.

The symplectic surface has its own quantization algebroid and its
fixed-class descent question.  Relating its centre to Hamiltonian
Lie cochains requires a comparison between those carriers.
Relating either centre to scalar traces requires an actual map to
the stated cyclic or representation-invariant target; surjectivity
is a further assertion.  The reduced trace homomorphism and the
stable Lie-homology theorem construct their stated maps without
making that additional identification.

